\documentclass[11pt]{article}
\usepackage[utf8]{inputenc}
\usepackage[T1]{fontenc}
\usepackage[spanish]{babel}
\usepackage{hyperref}
\usepackage[left=1.2cm,top=1.2cm,right=1.2cm,bottom=1.2cm]{geometry}
\usepackage{parskip}
\usepackage{enumitem}
\setlist[itemize]{topsep=1pt, itemsep=2pt, parsep=1pt, partopsep=0pt}
\usepackage{needspace}

\setlength{\parindent}{0pt}

\begin{document}

% =========================
% CABECERA (Apto para ATS)
% =========================
\begin{center}
    {\Large \textbf{Martín Alejandro Grellet}}\\[2pt]
    {\normalsize Senior Backend Engineer | Java · NestJS · Go}\\[2pt]
    Tucumán, Argentina · martin.alejandro.grellet@gmail.com ·
    \href{https://www.linkedin.com/in/mgrellet}{linkedin.com/in/mgrellet}
\end{center}

\vspace{6pt}

% =========================
% RESUMEN PROFESIONAL
% =========================
{\large \textbf{RESUMEN PROFESIONAL}}\\[4pt]
\hrule

Backend Engineer con más de 13 años diseñando sistemas distribuidos y APIs en Java (Spring Boot) y Node.js (NestJS / TypeScript). Experiencia en fintech, cumplimiento fiscal, e-commerce y logística, en entornos de alto volumen transaccional y producción crítica. He diseñado microservicios y pipelines orientados a eventos a escala, liderado decisiones técnicas y mentoreado equipos de ingeniería.

\vspace{10pt}

% =========================
% EXPERIENCIA PROFESIONAL
% =========================
{\large \textbf{EXPERIENCIA PROFESIONAL}}\\[4pt]
\hrule

\textbf{Siltium} \hfill Remoto, Argentina\\
\textit{Software Engineer – Correo Argentino} \hfill Mar 2025 -- Presente

\begin{itemize}[leftmargin=12pt]
    \item Diseño e implementación de una plataforma backend centralizada y de alto rendimiento en TypeScript y NestJS para registrar y gestionar datos de envíos internacionales, reduciendo la carga de trabajo manual en un $\sim$70\% y eliminando errores de entrada de datos.
    \item Construcción de una capa de API REST consumida por múltiples clientes internos, aplicando patrones de arquitectura limpia para garantizar la mantenibilidad a largo plazo de un equipo en crecimiento.
\end{itemize}

\textbf{AYI Group} \hfill Remoto, Argentina\\
\textit{Full Stack Software Engineer} \hfill Jul 2024 -- Ene 2025

\begin{itemize}[leftmargin=12pt]
    \item Liderazgo de la migración arquitectónica desde un monolito legado hacia microservicios (Java/Spring Boot) y microfrontends (React), mejorando la escalabilidad del sistema y reduciendo el riesgo en los despliegues.
    \item Optimización de planes de ejecución de consultas en MySQL y estrategias de indexación, disminuyendo los tiempos promedio de respuesta de la base de datos en un $\sim$35\% en momentos de alta carga.
\end{itemize}

\textbf{Sovos} \hfill Tucumán, Argentina\\
\textit{Principal Software Engineer} \hfill Ago 2022 -- Jul 2024

\begin{itemize}[leftmargin=12pt]
    \item Diseño e implementación de un motor de cumplimiento tributario en Java/Spring Boot para entornos POS, permitiendo cálculos impositivos precisos en tiempo real bajo complejas reglas transfronterizas.
    \item Reconstrucción de pipelines críticos de procesamiento por lotes en microservicios Java, reduciendo el tiempo de ejecución promedio y mejorando la respuesta de las APIs bajo alta concurrencia transaccional (de 1s a 100ms en promedio).
    \item Estandarización de flujos de trabajo de CI/CD utilizando Azure DevOps y Jenkins, eliminando pasos de lanzamiento manuales y reduciendo la fricción en los despliegues.
\end{itemize}

\textbf{Globant} \hfill Tucumán, Argentina\\
\textit{Java Software Designer – JP Morgan Chase} \hfill Mar 2022 -- Ago 2022

\begin{itemize}[leftmargin=12pt]
    \item Optimización de la gestión de dependencias de Maven y remediación de vulnerabilidades de seguridad identificadas por análisis estático con Veracode, reduciendo la exposición a CVEs críticos.
    \item Realización de pruebas, empaquetado y despliegue de microservicios Spring Boot en un entorno bancario regulado.
\end{itemize}

\textbf{Mercado Libre} \hfill Remoto, Argentina\\
\textit{Software Engineer} \hfill Sep 2021 -- Mar 2022

\begin{itemize}[leftmargin=12pt]
    \item Desarrollo de APIs REST en Go de alto rendimiento para servicios financieros y de pago (MercadoPago), construyendo endpoints que cumplen con normativas de Visa y procesando millones de transacciones diarias con tiempos de respuesta sub-100ms bajo picos de carga.
    \item Diseño de pipelines de pruebas de estrés automatizadas (Artillery) y establecimiento de tableros de observabilidad en tiempo real con Datadog y New Relic para asegurar el cumplimiento de SLAs.
\end{itemize}

\textbf{Sovos} \hfill Tucumán, Argentina\\
\textit{Software Engineer / Tech Lead} \hfill Feb 2016 -- Sep 2021

\begin{itemize}[leftmargin=12pt]
    \item Liderazgo de equipos multidisciplinarios de 4 a 6 ingenieros, impulsando decisiones técnicas, brindando mentoría y mejorando la previsibilidad en las entregas del equipo.
    \item Diseño y desarrollo de microservicios orientados a eventos para el cálculo de impuestos usando Java/Spring Boot y RabbitMQ, estableciendo la base de producción.
    \item Liderazgo de una prueba de concepto (PoC) entre Snowflake y Kafka para integrar datos en tiempo real, logrando incrementar el rendimiento de procesamiento en un $\sim$50\% en entornos de prueba.
\end{itemize}

\textit{Experiencia previa (2010--2016):} Java Developer / Application Support en Hewlett Packard Enterprise (servicios SOAP/REST de American Airlines, soporte 24/7) · Java Software Developer en Globant (portal interno de viajes con Java 7/Spring 3) · Desarrollador PHP en La Gaceta (portal digital de noticias) · Desarrollador de Software en Everis/HP (Java, Struts, Flex, PL/SQL)

\vspace{10pt}

% =========================
% EDUCACIÓN
% =========================
{\large \textbf{EDUCACIÓN}}\\[4pt]
\hrule

\textbf{Universidad del Norte Santo Tomás de Aquino} \hfill Tucumán, Argentina\\
Ingeniería en Informática

\vspace{10pt}

% =========================
% CERTIFICACIONES
% =========================
{\large \textbf{CERTIFICACIONES}}\\[4pt]
\hrule

\textbf{Oracle Certified Associate, Java SE 7 Programmer (1Z0-803)} \hfill 2014

\vspace{10pt}

% =========================
% HABILIDADES TÉCNICAS
% =========================
{\large \textbf{HABILIDADES TÉCNICAS}}\\[4pt]
\hrule

\textbf{Tecnologías principales:} Java (8, 11, 17), Spring Boot, Spring MVC, Hibernate/JPA, Node.js (NestJS), TypeORM, Go, TypeScript, JavaScript, SQL, Shell, React, AngularJS \\
\textbf{Bases de datos:} MySQL, PostgreSQL, SQL Server, Oracle, Snowflake \\
\textbf{Mensajería:} Kafka, RabbitMQ \\
\textbf{DevOps y herramientas:} Docker, Jenkins, Azure DevOps, Git, Maven, Gradle \\
\textbf{Pruebas (Testing):} JUnit, Mockito, Pruebas de Integración \\
\textbf{Prácticas:} Microservicios, REST APIs, CI/CD, DDD, SOLID

\vspace{10pt}

% =========================
% IDIOMAS
% =========================
{\large \textbf{IDIOMAS}}\\[4pt]
\hrule

\textbf{Inglés} (Competencia profesional) \quad \textbf{Español} (Nativo)

\end{document}
